% =====================================================================
\section{Walk Attention}\label{sec:model}
% =====================================================================

We now define the architecture. The guiding principle is the observation that the
multi-hop walk operator of Section~\ref{sec:algebra} is structurally an
\emph{attention}: it aggregates, for each target node, a weighted combination of
the features of the nodes that reach it. Transformer self-attention
\cite{vaswani2017attention} has the same shape but with all-pairs support and
\emph{learned} weights; graph attention \cite{velickovic2017graph} restricts the
support to one hop. Walk Attention keeps the learned weights of attention and
takes the support from the path algebra: the multi-hop, sparse walk-reachability
$\mathcal{S}_g$.

% ---------------------------------------------------------------------
\subsection{The Walk Attention layer}
% ---------------------------------------------------------------------

Fix a range $g$ with reachability support $\mathcal{S}_g$
(Definition~\ref{def:walkop}). A Walk Attention layer with $H$ heads and head
dimension $C$ transforms node features $x\in\RR^{n\times d_{\text{in}}}$ as
follows. For each head $h$ it computes queries, keys and values by linear maps,
\[
  Q^h = x\,W^h_Q,\qquad K^h = x\,W^h_K,\qquad V^h = x\,W^h_V,
  \qquad W^h_\bullet\in\RR^{d_{\text{in}}\times C},
\]
and then attends \emph{only over the reachable pairs}: for every
$(j,i)\in\mathcal{S}_g$,
\begin{align}
  \ell^h_{ji} &= \frac{1}{\sqrt{C}}\,\langle Q^h_j,\,K^h_i\rangle, \label{eq:logit}\\
  \alpha^h_{ji} &= \softmax_{\,i:\,(j,i)\in\mathcal{S}_g}\bigl(\ell^h_{ji}\bigr),
    \label{eq:softmax}\\
  z^h_j &= \sum_{i:\,(j,i)\in\mathcal{S}_g}\alpha^h_{ji}\,V^h_i. \label{eq:agg}
\end{align}
The softmax in \eqref{eq:softmax} is taken, for each target $j$, over exactly the
sources $i$ that reach $j$ by a length-$g$ walk---a \emph{segmented} softmax over
the support, never over all $n$ nodes. The head outputs are concatenated (or
averaged in the final layer) and a residual ``root'' term is added:
\[
  \mathrm{WA}_g(x)_j \;=\; \bigl\Vert_{h=1}^{H} z^h_j \;+\; x_j W_{\!\text{root}}.
\]
Because \eqref{eq:logit}--\eqref{eq:agg} only ever touch the pairs in
$\mathcal{S}_g$, the layer costs $O(|\mathcal{S}_g|\,H\,C)$ time and memory, not
$O(n^2)$. The support is precomputed once per graph from
Proposition~\ref{prop:count} (the nonzero pattern of $\Adj^{\,g}$).
Figure~\ref{fig:layer} summarises the data flow.

\begin{figure}[t]
\centering
\begin{tikzpicture}[
  >={Stealth[round]}, node distance=6mm,
  box/.style={draw,rounded corners=2pt,minimum height=7mm,inner sep=3pt,font=\small,fill=blue!4},
  alg/.style={draw,rounded corners=2pt,minimum height=7mm,inner sep=3pt,font=\small,fill=green!6},
  op/.style={font=\small},
  lbl/.style={font=\scriptsize,gray}]
  \node[box] (x) {$x$};
  \node[box,right=8mm of x] (q) {$Q=xW_Q$};
  \node[box,below=2mm of q] (k) {$K=xW_K$};
  \node[box,below=2mm of k] (v) {$V=xW_V$};
  \node[alg,right=10mm of k] (att)
    {$\displaystyle \alpha_{ji}=\softmax_{i:(j,i)\in\mathcal{S}_g}\!\frac{\langle Q_j,K_i\rangle}{\sqrt C}$};
  \node[box,right=10mm of att] (agg) {$z_j=\sum_i \alpha_{ji}V_i$};
  \node[box,right=8mm of agg] (out) {$\mathrm{WA}_g(x)$};
  % support callout
  \node[alg,below=8mm of att,fill=green!10] (supp)
    {support $\mathcal{S}_g=\{(j,i):(\Adj^{\,g})_{ij}>0\}$};
  \draw[->] (x) -- (q.west);
  \draw[->] (x) |- (k.west);
  \draw[->] (x) |- (v.west);
  \draw[->] (q.east) -- (att.north west);
  \draw[->] (k.east) -- (att.west);
  \draw[->] (att.east) -- (agg.west);
  \draw[->] (v.east) -- ([yshift=-2mm]agg.west);
  \draw[->] (agg) -- (out);
  \draw[->,green!50!black] (supp.north) -- (att.south)
    node[midway,right,lbl,text=green!50!black] {restricts pairs};
  \draw[->,rounded corners] (x.north) -- ++(0,6mm) -| (out.north)
    node[pos=0.25,above,lbl] {root (residual)};
\end{tikzpicture}
\caption{One Walk Attention layer at range $g$. Queries, keys and values are
linear maps of the node features; attention is computed \emph{only} over the
walk-reachable pairs $\mathcal{S}_g$ supplied by the path algebra (the nonzero
pattern of $\Adj^{\,g}$), then aggregated and added to a residual root term. The
green support box is what distinguishes Walk Attention from dense ($\mathcal{S}_g
=$ all pairs) and one-hop ($\mathcal{S}_g = $ edges) attention.}
\label{fig:layer}
\end{figure}

\begin{remark}[Why learned weights matter]
The fixed-weight operator $W^{(g)}$ of Definition~\ref{def:walkop} is the special
case in which $\alpha^h_{ji}$ is replaced by the constant $n_g(i,j)/Z_g$.
It reaches every range-$g$ source but weights them by multiplicity alone, so it
cannot \emph{select} which source carries the relevant signal. The learned
attention \eqref{eq:logit}--\eqref{eq:softmax} restores this selectivity while
keeping the same multi-hop support. Section~\ref{sec:exp} shows this distinction
is decisive.
\end{remark}

% ---------------------------------------------------------------------
\subsection{The network}
% ---------------------------------------------------------------------

A depth-$L$ Walk Attention network stacks one layer per range
$g=1,\dots,L$. Layer $g$ uses the support $\mathcal{S}_g$ derived from
$\Adj^{\,g}$, so successive layers read from progressively longer walks. Each
layer is followed by layer normalisation, an \textsc{elu} nonlinearity and
dropout; a final multilayer perceptron maps the node embedding to task outputs:
\[
  x^{(g)} \;=\; \mathrm{dropout}\Bigl(\mathrm{elu}\bigl(\mathrm{LN}\bigl(
  \mathrm{WA}_g\bigl(x^{(g-1)}\bigr)\bigr)\bigr)\Bigr),\qquad
  \hat y \;=\; \mathrm{MLP}\bigl(x^{(L)}\bigr).
\]
Layer normalisation is important in practice: without it the learned attention is
high-variance across initialisations (it reaches the optimum on some seeds and
not others); with it, together with a short learning-rate warmup and gradient
clipping, training converges reliably (Section~\ref{sec:exp}).
Algorithm~\ref{alg:wa} summarises one layer.

\begin{algorithm}[t]
\caption{Walk Attention layer at range $g$ (sparse).}\label{alg:wa}
\begin{algorithmic}[1]
\Require features $x\in\RR^{n\times d_{\text{in}}}$; support
  $\mathcal{S}_g=\{(j,i):(\Adj^{\,g})_{ij}>0\}$
\State $Q,K,V \gets xW_Q,\,xW_K,\,xW_V$ \Comment{per head; shapes $n\times HC$}
\For{each $(j,i)\in\mathcal{S}_g$ and head $h$}
  \State $\ell^h_{ji}\gets \langle Q^h_j,K^h_i\rangle/\sqrt{C}$
\EndFor
\State $\alpha^h_{ji}\gets$ segmented-$\softmax$ over $\{i:(j,i)\in\mathcal{S}_g\}$
\State $z^h_j\gets \sum_i \alpha^h_{ji}V^h_i$
  \Comment{scatter-add into targets}
\State \Return $\bigl\Vert_h z^h \;+\; xW_{\!\text{root}}$
\end{algorithmic}
\end{algorithm}

% ---------------------------------------------------------------------
\subsection{Relation to message passing and Transformers}
% ---------------------------------------------------------------------

Table~\ref{tab:relation} situates Walk Attention among related operators by two
axes: the \emph{support} of interaction and whether the weights are
\emph{learned}. A one-hop graph-attention layer \cite{velickovic2017graph} learns
weights but only over direct neighbours, so it must propagate long-range signal
through $L$ stacked layers---precisely the regime where over-squashing is most
severe. A full graph Transformer \cite{vaswani2017attention} learns weights over
all pairs,
paying $O(n^2)$ and discarding graph structure unless it is reinjected by
positional encodings. The fixed-weight walk operator has the right multi-hop
support but cannot select. Walk Attention occupies the remaining cell: multi-hop,
structure-aware support \emph{and} learned, selective weights.

\begin{table}[t]
\centering
\caption{Walk Attention against related aggregation operators.}
\label{tab:relation}
\begin{tabular}{lccc}
\toprule
\textbf{Operator} & \textbf{Support} & \textbf{Weights} & \textbf{Cost} \\
\midrule
GCN / message passing       & 1-hop      & fixed (degree) & $O(|Q_1|)$ \\
Graph attention (GAT)       & 1-hop      & learned        & $O(|Q_1|)$ \\
Walk operator $W^{(g)}$     & walk-reachable & fixed (multiplicity) & $O(|\mathcal{S}_g|)$ \\
Graph Transformer           & all pairs  & learned        & $O(n^2)$ \\
\textbf{Walk Attention (ours)} & \textbf{walk-reachable} & \textbf{learned} & $O(|\mathcal{S}_g|)$ \\
\bottomrule
\end{tabular}
\end{table}
